\section*{Chapter \ref{ch:s1:flows} --- Flows}

\begin{solution}{exo:s1:flows:1}
$0.2 \times 0.05 = 1.0\%$ of the stock's shares to buy; with an impact of 1.5, a push of 1.5\% over the quarter.
\end{solution}

\begin{solution}{exo:s1:flows:2}
Half (50\%) after 126 days, a quarter after 252: three quarters of the push has reversed within a year.
\end{solution}

\begin{solution}{exo:s1:flows:3}
Filed on the last day allowed, the March 31 holdings are 45 days old when they arrive (mid-May); they remain the latest until the June holdings arrive in
mid-August, when they are about four and a half months old.
\end{solution}

\begin{solution}{exo:s1:flows:4}
Funds that did well hold stocks that did well; the synthetic stocks have persistent drifts, so the stocks the good funds hold keep rising. Flows chase the
funds' returns and are traded in those same stocks, so \omterm{def:s1:flows:fit}{flow-induced trading} and past returns line up even with no pressure (a cross-sectional slope of 0.31).
\end{solution}

\begin{solution}{exo:s1:flows:5}
A fund with good returns gets inflows, buys more of what it holds, and pushes those prices up; that raises its next returns, which attract more inflows.
Its performance persists because of the flows it attracts, not its skill: Lou's explanation.
\end{solution}

\begin{solution}{exo:s1:flows:6}
The trade can only start when the quarter's trading is known, 45 days after the quarter, by which time part of the reversal has happened; what remains is
small and is fought by the drifts that the flows were chasing, which keep the sold stocks falling. The \omterm{def:s1:earnings:event}{event study} measures the reversal from the quarter's
start; the book enters much later.
\end{solution}

\begin{solution}{exo:s1:flows:7}
The top-minus-bottom decile gains 4.9\% in the quarter and 69.0\% over three years. Funds tilted to a rewarded style earn its premium, attract inflows and
buy that style, so the flow signal ranks stocks by their style. To test a real flow signal, regress its returns on the style factors and momentum, or
build it from flows residualised on the fund's style returns.
\end{solution}

\begin{solution}{exo:s1:flows:8}
Quarter-end holdings are filed up to 45 days later, so trading on the quarter-end date uses information not yet public; flows are also reported with a lag.
The honest backtest trades on the filing dates, and controls for momentum and styles.
\end{solution}

\begin{solution}{pb:s1:flows:1}
\begin{enumerate}
\item Trading forced by fund flows (ownership share times flow, summed over funds); the part of a price move caused by uninformed demand, expected to reverse.
\item Funds with large outflows sell existing positions, creating pressure in commonly held stocks; trading against them earns significant returns; the
forced trades are predictable.
\item A significant, temporary price impact of \omterm{def:s1:flows:fit}{flow-induced trading}; its expected part forecasts next year's returns, reversed later; it can explain fund
performance persistence and the smart-money effect, and part of momentum.
\item Institutions are constrained, the market's demand is inelastic, and \$1 invested raises its value by about \$5.
\item Two hundred funds of sixty random stocks, a quarter replaced each quarter, owning 20\% of each stock; flows of 5\% per standard deviation of trailing
return plus 3\% noise; a push of 1.5 per unit of \omterm{def:s1:flows:fit}{flow-induced trading} over the quarter, reversing with a half-life of 126 days.
\item The push raises the returns of the funds that hold the pushed stocks, which then attract more flows.
\item 1.05 in the quarter and $-0.69$ the next year with pressure; 0.31 and 0.32 without.
\item 66\% (75\% of the end-of-quarter push by the planted decay alone).
\item Flows forecast from trailing fund returns times filed holdings; Sharpe ratio 0.44 after costs.
\item It enters 45 days after the quarter, too late for much of the reversal, and fights the drifts; $-0.35$.
\item Mechanical basket trades through arbitrage: more volatility and negative autocorrelation in the stocks they hold.
\item Quarterly holdings reports of managers above \$100 million in 13(f) securities, filed within 45 days.
\item \textbf{Named result.} Each unit of \omterm{def:s1:flows:fit}{flow-induced trading} moves the quarter's log return by 1.05 (0.31 of it the planted drift the flows chase), and 66\%
of the quarter's move reverses within the following year.
\item The top-minus-bottom decile gains 1.74\% over the quarter, peaks at 1.91\% just after it, is back to zero after about a year and ends at $-2.15\%$ after
three.
\item With style-tilted funds the signal becomes a style bet: 69.0\% over three years, never reversing.
\item Monthly fund assets and flows, daily ETF creations, and fund-level reports of portfolio changes.
\item Regress its returns on momentum and the style factors, or residualise the flows on fund style returns.
\item The expected-flow book, with the flow model re-estimated as flows change.
\item Faster knowledge of the forced selling: daily flow data or news of liquidations.
\item Liquidity, to investors who must trade.
\end{enumerate}
\end{solution}

\begin{solution}{iq:s1:flows:1}
Because funds must trade their flows in the stocks they hold, and the other side of that demand is not perfectly elastic; the price moves until someone is
paid enough to take it, and moves back as the demand is absorbed.
\end{solution}

\begin{solution}{iq:s1:flows:2}
Holdings from 13F and fund portfolio filings, flows from changes in fund assets net of returns, each dated by publication; \omterm{def:s1:flows:fit}{flow-induced trading} is the sum over
funds of ownership share times flow. The lag is weeks to months: holdings up to 45 days after the quarter for 13F, flows monthly or quarterly.
\end{solution}

\begin{solution}{iq:s1:flows:3}
Regress its long--short returns on momentum and style factors, sort within momentum deciles, and build the flow measure from flows residualised on fund
style and past returns; a signal that survives is about flows.
\end{solution}

\begin{solution}{iq:s1:flows:4}
Estimate what it must sell and how fast, trade against it only once its selling has pushed prices below value, and hold for the reversal. Worry about other
front-runners, the fund selling the same names as others under stress, and whether the liquidation reflects information.
\end{solution}

\begin{solution}{iq:s1:flows:5}
Correlated selling: when those funds face outflows, the stock is sold with its other common holdings regardless of its own news, so its risk includes the
funds' flow risk.
\end{solution}

\begin{solution}{iq:s1:flows:6}
The remaining push $s$ days after the quarter is $\lambda F\,2^{-s/h}$. Buying (for $F < 0$) at $L$ and selling at $L + H$ earns
$-\lambda F\,2^{-L/h}(1 - 2^{-H/h})$, positive for sold stocks, before costs $2c$; the trade pays when $|\lambda F|\,2^{-L/h}(1 - 2^{-H/h}) > 2c$. A long lag
$L$ shrinks it exponentially, which is why the fire-sale book needs fast data.
\end{solution}
